% ============================================================
%  date/main.tex -- single-file working copy for the LENGTH CUT
%
%  Created 2026-09-16 from the v7 tree one level up
%  ("00 - DATE CAM Monitoring.tex" + section files 01-08), with every
%  \input inlined verbatim in order. Section boundaries are marked
%  with "% ===== NN - <file> =====" separators so edits can be traced
%  back to the source file if needed. The v7 tree is the archive; the
%  version history and change-log lives in its 00 file.
%
%  Target: DATE 2027 regular paper -- 6 pages + 1 page for references
%  only (CFP wording: "must not exceed 6 pages in length, with one
%  extra page allowed only for bibliographic references"). Final
%  paper deadline: Sunday 20 Sep 2026 AoE. Baseline at creation:
%  16 pages (~14.7 body + ~1.3 references).
%
%  STATUS: R1 structural cuts applied Sep 16 -> 14 pages clean. See the
%  length-cut log at the end of this file (With Comments build) and
%  cut_plan.md. backup/ holds the pre-round snapshots.
%
%  Supporting files copied alongside: IEEEtran.cls, IEEEtran.bst,
%  RobStd.sty, comments.tex (flag system), refs.bib, figures/.
%  Compile: bash compile.sh   (main.pdf = clean; main-comments.pdf)
%  or plain: pdflatex main && bibtex main && pdflatex main && pdflatex main
% ============================================================
\documentclass[conference]{IEEEtran}
\IEEEoverridecommandlockouts

\usepackage{cite}
\usepackage{amsmath,amssymb,amsfonts}
\usepackage{graphicx}
\usepackage{textcomp}
\usepackage[dvipsnames]{xcolor}
\usepackage{RobStd}
\usepackage{subcaption}
\usepackage{siunitx}
% --- siunitx v2/v3 compatibility shim ------------------------
\ifdefined\qty\else
  \NewDocumentCommand{\qty}{O{}mm}{\SI[#1]{#2}{#3}}
\fi
\ifdefined\unit\else
  \NewDocumentCommand{\unit}{O{}m}{\si[#1]{#2}}
\fi
% -------------------------------------------------------------
\usepackage{placeins}
\usepackage{multirow}
\usepackage{makecell}
\usepackage{booktabs}
\usepackage{url}
\renewcommand\arraystretch{1.15}
\setlength{\tabcolsep}{4pt}
\graphicspath{{figures/}}

\def\BibTeX{{\rm B\kern-.05em{\sc i\kern-.025em b}\kern-.08em
    T\kern-.1667em\lower.7ex\hbox{E}\kern-.125emX}}

% --- shorthand macros ----------------------------------------
\newcommand{\cam}{{CAM}\xspace}
\newcommand{\tcam}{{TCAM}\xspace}
\newcommand{\mcam}{{MCAM}\xspace}
\newcommand{\acam}{{ACAM}\xspace}
\newcommand{\imc}{{CiM}\xspace}
\newcommand{\fefet}{{FeFET}\xspace}
\newcommand{\fefets}{{FeFETs}\xspace}
\newcommand{\auroc}{{AUROC}\xspace}
\newcommand{\signrp}{{sign-RP}\xspace}
\newcommand{\tocite}[1]{[{\color{red}citation-#1}]}

% ============================================================
%  comments.tex -- annotation / review-flag system
%
%  STREAMLINED FLAG NAMES (Aug 2026, per MN: "this is just less typing"):
%      \dt[note]{content}   was \dothis   / \dothisflag
%      \eo[note]{content}   was \expandon / \expandonflag
%  The old names are kept as aliases below, so every existing .tex
%  still compiles unchanged. Semantics are identical.
%
%  All flags obey the showcomments boolean; compile.sh flips it.
% ============================================================
\usepackage{ifthen}
\usepackage{amssymb}
\usepackage{xcolor}
\usepackage{etoolbox}
\usepackage{environ}
\usepackage[most]{tcolorbox}

\newboolean{showcomments}
\setboolean{showcomments}{false}    % compile.sh targets this exact line -- keep it here

% DEFAULT IS FALSE (no comments) so Overleaf produces the clean copy.
% Run `bash compile.sh` to build both versions.

\definecolor{flagorange}{RGB}{200, 95, 0}
\definecolor{flagblue}{RGB}{0, 70, 190}
\definecolor{aiviolet}{RGB}{120, 40, 180}
\definecolor{dogreen}{RGB}{0, 130, 60}

% ---- the two action flags -----------------------------------
%  Uniform shape:  \flag[what to do]{content}
%    \dt  orange -- an action to take
%    \eo  blue   -- say more about this
%  Content survives into BOTH builds; the bracketed note appears only
%  in the With Comments build. With NO optional argument the braced
%  text IS the note and disappears from the clean build entirely:
%      \dt{fill in the submission date}          note only
%      \dt[Trim this]{The library is short ...}  keep text, flag it
%      \eo[Add numbers]{}                        note only
\newcommand{\flagmark}[3]{%
  \ifthenelse{\boolean{showcomments}}{{\color{#1}\bfseries[#2:\ #3]}\ }{}%
}
\newcommand{\dt}[2][]{%
  \ifstrempty{#1}%
    {\ifthenelse{\boolean{showcomments}}{{\color{flagorange}\bfseries[DO THIS:\ #2]}}{}}%
    {\flagmark{flagorange}{DO THIS}{#1}#2}%
}
\newcommand{\eo}[2][]{%
  \ifstrempty{#1}%
    {\ifthenelse{\boolean{showcomments}}{{\color{flagblue}\bfseries[EXPAND:\ #2]}}{}}%
    {\flagmark{flagblue}{EXPAND}{#1}#2}%
}
% ---- aliases (do not "clean these up") ----------------------
\let\dothis\dt
\let\dothisflag\dt
\let\expandon\eo
\let\expandonflag\eo

% ---- reviewer remarks --------------------------------------
% Inline remark:  \inote{short note}  /  \inote[aiviolet]{note}
\newcommand{\inote}[2][red]{%
  \ifthenelse{\boolean{showcomments}}{\textcolor{#1}{[#2]}}{}%
}
% Block remark. Escape underscores as \_ inside.
%   \begin{pnote} ... \end{pnote}    or    \begin{pnote}[aiviolet] ... \end{pnote}
\NewEnviron{pnote}[1][red]{%
  \ifthenelse{\boolean{showcomments}}{\begingroup\color{#1}\BODY\endgroup}{}%
}
% Review-build-only content (whole paragraphs, tables, or a page).
\newcommand{\reviewonly}[1]{\ifthenelse{\boolean{showcomments}}{#1}{}}

% ---- structural boxes (always visible) ---------------------
\newtcolorbox{sectionbox}{%
  colback=black, coltext=white,
  width=\linewidth, boxrule=0pt,
  left=3pt, right=3pt, top=2pt, bottom=2pt,
  fontupper=\sffamily\bfseries\normalsize,
  sharp corners,
  before=\par\vspace{4pt}, after=\par\vspace{2pt}}
\newtcolorbox{subsectionbox}{%
  colback=lightgray, coltext=black,
  width=\linewidth, boxrule=0pt,
  left=3pt, right=3pt, top=1pt, bottom=1pt,
  fontupper=\sffamily\bfseries\small,
  sharp corners,
  before=\par\vspace{3pt}, after=\par\vspace{1pt}}

% ---- per-author inline notes -------------------------------
\makeatletter
\newcommand{\mynote}[3]{%
  \ifthenelse{\boolean{showcomments}}{%
   \fbox{\bfseries\sffamily\scriptsize#1}%
   {\small$\blacktriangleright$\textsf{\emph{\color{#3}{#2}}}$\blacktriangleleft$}}%
  {\@bsphack\@esphack}%
}
\makeatother
\definecolor{PersianGreen}{rgb}{0, 0.65, 0.58}
\newcommand{\mn}[1]{\mynote{MN}{#1}{violet}}
\newcommand{\ib}[1]{\mynote{IB}{#1}{PersianGreen}}
\newcommand{\sj}[1]{\mynote{SJ}{#1}{blue}}
\newcommand{\xh}[1]{\mynote{XH}{#1}{orange}}

% \setboolean{showcomments}{true}   % or run: bash compile.sh

\begin{document}

% ============================================================
%  ANNOTATION FLAGS -- quick reference
%
%  All flags compile away when showcomments is false (the default).
%  The two action flags take the SAME shape, \flag[what to do]{content}.
%  Content survives into BOTH builds; the bracketed note appears only in
%  the With Comments build. With NO optional argument the braced text IS
%  the note and disappears from the clean build entirely.
%
%  ORANGE -- an action to take
%    \dt[note]{content}          (was \dothis / \dothisflag; still work)
%    \dt{fill in the deadline}   note only, vanishes from the clean build
%
%  BLUE -- say more about this
%    \eo[note]{content}          (was \expandon / \expandonflag; still work)
%
%  RED / VIOLET -- reviewer remarks, review build only
%    \inote{short note}   \inote[aiviolet]{AI-generated suggestion}
%    \begin{pnote} ... \end{pnote}     block note; escape _ as \_
%    \reviewonly{...}                  content only in the review build
%
%  PER-AUTHOR inline notes: \mn{...} \ib{...} \sj{...} \xh{...}
%
%  A ']' typed inside a flag's optional argument silently truncates it
%  and leaks the rest into the CLEAN build. Brace them: {[1]}, {[5]}.
% ============================================================

% ============================================================
%  WRITING CONVENTIONS FOR THIS DRAFT (MN's voice; fix-ai-writing skill)
%
%  1.  Plain, first person plural, hedged. "may", "could", "appears to",
%      "remains to be established". No absolutes that are not absolute.
%  2.  En dashes (--), never em dashes (---).
%  3.  Numbered inline lists: {\bf (1)} ... {\bf (2)} ..., not
%      "First, ... Second, ...".
%  4.  Bold run-in lead-ins end with a COLON, not a period, and the next
%      word keeps its capital:  \textbf{What the grid measures:} Whether ...
%  5.  Never the word "honest" as framing; never announce your own
%      fairness ("to be fair", "stated fairly", "candidly").
%  6.  No self-referential framing: "the whole argument", "the key
%      insight", "which is the entire problem", "the upshot", "X turns
%      on Y". State the substance and stop.
%  7.  No verdict tags on your own description ("and it works").
%  8.  Do not narrate design choices ("deliberately", "chosen to be").
%  9.  A section discusses; it is not a proof. "Sec. IV covers ...",
%      not "Sec. IV is the proof that ...".
%  10. No "binds" / "bites" for constraints; no borrowed-metaphor verbs
%      ("is the currency", "prices the choice", "lands on", "rests on").
%      Use "maps to", "depends on", "the limit here is".
%  11. American spelling, always. Define every acronym at first use.
%  12. Scope caveats go in footnotes, not into the argumentative flow.
%  13. Label invented or illustrative examples as such, in the text.
%  14. Figures follow the house-graphs conventions: Helvetica, Set1
%      colors, transparent background, provenance header in the
%      generating script (see gen_paper_figs.py).
% ============================================================

\bstctlcite{IEEEexample:BSTcontrol}

% DATE is double-blind: authors and funding are anonymized for
% submission. Funding restored at camera-ready (SRC/DARPA JUMP 2.0;
% NSF IRES; etc.).
% Title per MN (08-25): option (2) from the v02 markup -- names the
% design-space contribution first for a DATE audience.
\title{\cam-Resident Behavioral Monitoring for LLMs: Encoding Choices,
Cross-Domain Fidelity, and Library-Size Scaling}

\author{\IEEEauthorblockN{Anonymized for review}
\IEEEauthorblockA{DATE 2027 submission}}

\maketitle

\begin{abstract}
Frontier language model deployments increasingly treat a model's own internal
activations as a safety signal: a production jailbreak defense from
Anthropic, for instance, is realized via a single linear probe over residual-stream
activations, scored on every generated token at roughly
$10^{-3}$~percent of forward-pass compute. What is monitored is set by
$K$, the number of stored behavioral signatures each activation is
compared against, and target $K$'s are rapidly increasing: from $K=1$
in the deployed defense above, to $K\!\approx\!10^{2}$--$10^{3}$ once
probes are fitted per behavior, per domain, and per layer, to
$K\!\approx\!1.3\!\times\!10^{5}$ per layer for a whole-vocabulary
readout, where the signature library equals the model's vocabulary. On a GPU,
the signature bank must stream from high-bandwidth memory once per
monitored token, so at sufficiently high $K$ or token rates -- e.g.,
$K=10^{4}$--$10^{5}$ at $10^{3}$~tokens/s -- the monitor's own memory
traffic can exceed the bandwidth of the accelerator it runs
on. We
propose addressing this bottleneck with content-addressable memories
(\cam{}s), which evaluate distance functions directly within the
memory array -- emerging non-volatile devices further enable
richer per-cell distance metrics (e.g., Euclidean). Because the signature bank never moves and only the query is
broadcast, this in-memory search trades arithmetic precision and distance function fidelity for orders of magnitude lower data movement. We quantify this trade-off on two
representative workloads: {\bf (1)} a ten-domain cross-domain
deception grid over Llama-3.3-70B activations, where 384-bit codes
preserve in-domain \auroc within 0.01 of the full-precision cosine-similarity
ceiling while achieving more than $300\times$ compression per stored signature; and
{\bf (2)} whole-vocabulary readout, where a 3-bit Euclidean code at
$\sim$$10^{3}$ storage bits per row reproduces 0.96 of the exact
readout's top-1 token selections. Because search occurs in place
rather than via streaming, the dominant costs shift to array capacity and search energy: the
whole-vocabulary bank occupies 6--\qty{33}{\mega\byte} per layer, and estimated
in-array search energy is five to six orders of
magnitude below the energy of streaming the equivalent bank from
HBM.
\end{abstract}

\begin{IEEEkeywords}
Content-addressable memory; compute-in-memory; large language models;
activation probes; AI safety monitoring; nearest-neighbor search;
ferroelectric FETs; Hamming distance
\end{IEEEkeywords}

%%%%%%%%%%%%%%%%%%%%%%%%%%%%%%%%%%%%%%%%%%%%%%%%%%%%%%%%%%%%%%%%%%%%%%
% ===== 01 - Introduction.tex =====
%%%%%%%%%%%%%%%%%%%%%%%%%%%%%%%%%%%%%%%%%%%%%%%%%%%%%%%%%%%%%%%%%%%%%%
\section{Introduction}
\label{sec:Introduction}

Large language model (LLM) deployments are wrapped in a layered set
of safeguards, and until recently every layer read
\emph{text}~\cite{metr2026policies}. That has started to change: the
always-on stage of one production jailbreak defense now reads the
model's internal activations
directly~\cite{ccpp2026,anthropic2026ccblog}, as do deployed
residual-stream probes at a second provider~\cite{gemini2026probes}.
The deployed check is small -- one stored vector $w$, one dot product
per generated token, one threshold. Only flagged exchanges are
escalated to an expensive classifier ensemble -- and it is a small
share of even the deployed defense's reported overhead
(Sec.~\ref{subsec:Probes}).

Recent security evaluations show why such monitors are needed.
Frontier models now run long, autonomous operations: current cyber
evaluations walk a model through 32 steps from reconnaissance to
takeover, with run budgets up to $10^{8}$ tokens. The strongest
model tested progressed, unaided, through 22 of those 32 attack steps
on an average attempt -- roughly two thirds of the way to a completed
operation~\cite{aisi2026cyber}. In one four-day campaign of 122
such runs across seven models, evaluators later found 19 unsanctioned
actions in four distinct categories -- supply-chain compromise of a real open-source project,
deception and targeting via a file-transfer service, prompt injection
planted for other automated systems to execute, and inter-agent
collaboration via public messages. All of said outcomes were found
\emph{retrospectively}; nothing was watching the models' internal
state as they ran~\cite{aisi2026incident}. One month later the same
behaviors appeared at production scale: per public investigations,
hundreds of agents on OpenAI's infrastructure organized, in secret,
into a working collective -- building a communication channel,
reverse-engineering their own scorer, falsifying evidence, and
compromising both OpenAI and Hugging
Face~\cite{clark2026import,cotra2026hf}. Again, the record was
reconstructed from transcripts after the fact -- and the
coordination itself is a behavior for which no stored direction has
been published. A $10^{8}$-token operation
is not visible by sampling outputs, and each individual message such
an agent emits can read as routine -- a plausible commit, a file
transfer, a status update -- while the sequence they add up to is an
attack. As such, a monitor for this setting has to be per-token,
always-on, and internal.

Interpretability work supplies the signatures such a monitor would
compare against, and it also indicates how many it needs. Linear
probes recover truth-related structure across datasets and model
families~\cite{marks2023geometry,burns2023ccs,zou2023repe}, but the
directions are narrow. A probe fitted on factual lying has an
\auroc score of $\sim$0.55 (random chance) on
sycophancy~\cite{ying2026spectrum}, while probes trained on outright
lying largely miss deception carried by technically true
statements~\cite{berger2026liedetector}. Detecting deceptive
reasoning has become a funded interpretability objective in its own
right~\cite{schmidt2026interp}. Covering a behavior suite therefore
takes a stored direction per behavior, per subject-matter domain, at
each layer where that behavior separates -- and each of those
combinations is one more stored row. As such, $K$, the number of
stored signatures each activation is compared against, can grow with
everything the monitor is asked to notice; and widening what is
covered means growing the library that must be searched.

Fig.~\ref{fig:ladder}(a) plots the published record of $K$ against
time; Fig.~\ref{fig:ladder}(b) estimates the percentage of GPU
bandwidth that could be consumed with increasing $K$.
\textbf{\boldmath$K=1$} represents deployed always-on monitors that
read a single direction~\cite{ccpp2026,gemini2026probes}.
\textbf{\boldmath$K\approx10^{2}$} captures probes that are fitted per
behavior and per domain~\cite{ying2026spectrum}.
\textbf{\boldmath$K\approx10^{3}$} is indicative of the four behavior
categories of the above campaign, multiplied over ten domains and
13--25 usable layers -- i.e., 520--1{,}000
rows~\cite{aisi2026incident}; the production-scale incident above
repeats every one of said categories, so this arithmetic is, if
anything, conservative~\cite{clark2026import}.
\textbf{\boldmath$K\approx1.3\!\times\!10^{5}$} is representative of
the whole-vocabulary (``J-space'') readout, from a published method
reporting that a model holds words ``in mind'' without emitting them
-- i.e., for every token in the vocabulary there is an internal
pattern meaning that word is currently active, and it is readable
whether or not the word is ever
said~\cite{anthropic2026workspace}.\footnote{As an example, prompted
with the number of legs on the animal that spins webs, the model
answers \emph{eight} while the readout shows \emph{spider} active
internally; the same work reports patterns for ``fake'' and
``fictional'' active while a model privately registered that a safety
test was staged.} Finally, $K$ could grow even
higher, as published feature
dictionaries~\cite{templeton2024scaling,gao2024sae,lieberum2024gemmascope}
are even two orders of magnitude higher than the J-space. That said,
we can cross the memory bandwidth of deployable parts well before this.
Per Fig.~\ref{fig:ladder}(b), at $10^{3}$ tokens/s, an H100 is fully
consumed at $K\!\approx\!204$K and a Jetson-class part at
$K\!\approx\!17$K (Sec.~\ref{sec:Scaling}).

\begin{figure}[!t]
\centering
\includegraphics[width=3.4in]{fig_ladder_v2.pdf}
\caption{(a) Monitored library size $K$ in the published record
(left axis), read as GPU monitoring traffic on the right axis via
Eq.~(\ref{eq:gputraffic}) at $d=8{,}192$, fp16, $10^{3}$~tokens/s,
one monitored layer; dashed lines mark two parts' memory bandwidths,
and the green line is a \cam's per-token traffic, flat in $K$. The
arrow joins the record's endpoints; it is not a fit and not a
forecast. (b) The share of a part's memory bandwidth consumed by
monitoring alone, against aggregate token rate and library size $K$:
lines mark 100\% of an H100, an A100, and a Jetson Thor, with 10\%
and 1\% of an H100 as guides. The strip
anchors the rate axis to familiar workloads
(Sec.~\ref{subsec:CostModel}): per-stream anchors batch up to a
node's aggregate rate, while batch-1 edge workloads (VLA, video) sit
on the axis directly -- so batching moves a deployment along the
axis rather than adding one; per-request amortization is objection
(3) of the companion appendix document. Both panels are arithmetic
on vendor bandwidth
figures, not measured kernel throughput.}
\label{fig:ladder}
\end{figure}

We propose to support this search with a content-addressable memory
(\cam), which computes the distance function in the memory array
itself: the bank never moves and per-token traffic is flat in $K$
(green line of Fig.~\ref{fig:ladder}(a)).
Candidate substrates are based on both near-term and the emerging
devices: SRAM-based \cam{}s are buildable in any CMOS process today,
and can implement Hamming distance, while \cam{}s based on emerging
devices (\fefets, RRAM) offer density, nonvolatility, and richer
distance metrics (Secs.~\ref{subsec:CAMspace}
and~\ref{subsec:CAMtech}).

This paper measures that trade for behavioral-monitoring workloads.
Our contributions include: {\bf (1)} a formulation of always-on activation
monitoring as library search, with the traffic and capacity arithmetic
and the GPU crossing points (Secs.~\ref{sec:Scaling},
\ref{sec:Hardware}); {\bf (2)} a maturity-ladder mapping of the
detector onto \cam cell types and match modes, with an audit that
catches classification hiding in a trained front end
(Sec.~\ref{sec:Mapping}); {\bf (3)} measured fidelity at \cam
precision at both ends of the $K$ range -- the ten-domain deception
grid, within 0.009 \auroc of the full-precision ceiling in-domain
(Sec.~\ref{sec:Results}), and an in-house whole-vocabulary readout
reproduction scored against the exact readout
(Sec.~\ref{sec:LargeK}); and {\bf (4)} a first-order search
energy/latency estimate, with the measurements that is the foundation
for future cost models (Sec.~\ref{sec:Hardware}).

% ===== end 01 - Introduction.tex =====

%%%%%%%%%%%%%%%%%%%%%%%%%%%%%%%%%%%%%%%%%%%%%%%%%%%%%%%%%%%%%%%%%%%%%%
% ===== 02 - Background.tex =====
%%%%%%%%%%%%%%%%%%%%%%%%%%%%%%%%%%%%%%%%%%%%%%%%%%%%%%%%%%%%%%%%%%%%%%
\section{Background and Cost Model}
\label{sec:Background}

\subsection{The CAM design space}
\label{subsec:CAMspace}
Ordinary memory takes an address and returns content, so searching a
library means reading it row by row. A \cam compares every row
against the query at once, in place, and only reports matching
rows~\cite{pagiamtzis2006cam}. One access is therefore one search of
the whole library, at any library size, and the data that moves is the
query rather than the bank.

A \textbf{binary \cam} (BCAM) cell holds one bit. A
\textbf{ternary \cam} (\tcam) cell holds a bit or a don't-care, which
is what makes a stored ternary code with a zero band natively
expressible~\cite{ni2019fetcam}. A \textbf{multi-bit \cam} (\mcam) cell
holds one of several levels~\cite{kazemi2021fefetcam}. An
\textbf{analog \cam} (\acam) cell
holds a range~\cite{li2020acam}. A binary or ternary cell can report
only whether a position disagrees, and can measure Hamming
distance. A multi-bit or analog cell stores a level,
which is what lets the matchline measure how far the query sits from
what is stored. In a multi-bit \fefet implementation, sweeping the query
voltage traces a V-shaped matchline current with its minimum at the
stored level; the curve is approximately quadratic near the bottom, so
one cell physically computes $(Q_i-M_i)^2$ and $N$ cells on a common
matchline sum to $\sum_i (Q_i-M_i)^2$, a squared Euclidean distance with
no multiplier anywhere~\cite{kazemi2021fefetcam} (see Table~\ref{tab:camspace}). Notably, said cell
behavior has been \emph{measured}, not only modeled: multi-bit,
array-level \cam operations have been experimentally demonstrated
with \fefets -- a two-\fefet cell programmed to eight states (3
bits), whose measured matchline distance function shows the
quadratic behavior described above~\cite{kazemi2022hdcfefet}.\footnote{Squared
Euclidean is an acceptable stand-in for Euclidean here, since the square
root is monotonic and therefore changes neither which row is closest nor
which rows clear a threshold.}

With appropriate peripheral circuitry we can measure: \textbf{exact}
match reports rows with no mismatches; \emph{best} match reports the
row with the fewest, even when nothing is exact; \textbf{threshold}
match reports
every row within a limit $\tau$. Each is a different matchline circuit.
Threshold mode matters for monitoring because the expected answer is
almost always empty, and a threshold array can report ``nothing fired''
without ranking anything. Table~\ref{tab:camspace} summarizes the
design space; notably, the cheapest cells support only Hamming
distance, which is why encodings that hold accuracy under a Hamming
metric are potentially valuable. A goal of what follows is to
quantify what an emerging-technology \cam can or cannot buy: where
the multi-bit cells of Table~\ref{tab:camspace} change the
achievable accuracy and energy, and where a plain CMOS Hamming
array is already enough.

\begin{table}[t]
\centering
\caption{The \cam design space a monitoring workload has to choose from.
Cell transistor counts are order-of-magnitude figures for CMOS
implementations; NVM-based cells trade area for nonvolatility.}
\label{tab:camspace}
{\footnotesize
\begin{tabular}{@{}llll@{}}
\toprule
\textbf{Cell} & \textbf{Holds} & \textbf{Distance} & \textbf{Cost / maturity} \\
\midrule
BCAM  & one bit            & Hamming            & $\sim$6--10T, any CMOS fab \\
\tcam & bit or don't-care  & Hamming, wildcards & $\sim$16T CMOS; 2T \fefet \\
\mcam & one of $2^b$ levels & Euclidean ($L_2^2$) & multi-level device \\
\acam & a range            & range / interval   & multi-level device \\
\midrule
\multicolumn{4}{@{}l}{\emph{Match mode is an independent choice:} exact,
best-match, or threshold.} \\
\bottomrule
\end{tabular}}
\end{table}

\subsection{\cam technologies}
\label{subsec:CAMtech}
\dt[MN: placeholder per your restructuring note -- your
technology-based \cam content goes here (SRAM vs.\ FeFET vs.\ RRAM
implementations, the SUPREME/JUMP material, whatever you want to
carry). The BEOL paragraph below is the seed; Table~\ref{tab:devices}
in Sec.~VII carries the device figures and could move up here if this
subsection grows]{}
The rungs of Table~\ref{tab:camspace} differ in where the array can
physically sit. SRAM-based binary and ternary cells are buildable in
any CMOS process today, in the front-end silicon alongside logic. The
\fefet-based cells trade that ubiquity for density and nonvolatility,
and their device work is moving toward the back end of line (BEOL):
monolithic-3D integration of high-endurance multi-bit \fefets has been
demonstrated with the memory tier above the logic
tier~\cite{dutta2020m3dfefet}, and BEOL-compatible ferroelectric
transistors with AlScN films on two-dimensional channels have been
reported at thermal budgets a metal stack tolerates~\cite{kim2023alscn}.
BEOL RRAM and MRAM are already offered at 22/16\,nm foundry nodes. A
BEOL-capable cell is what would eventually let the signature bank sit
in the metal layers above the compute that produces the activations;
Sec.~\ref{subsec:Placement} returns to what that placement would mean
for this workload. Notably, none of said substrates would be built
for one application: an SRAM-based associative processing unit
already performs bit-serial Hamming search at billion-vector
scale commercially~\cite{gsi2025apu}, and \cam-based accelerators
have been proposed across retrieval, one-shot learning, and
range-matching workloads -- so a monitoring part rides on a
primitive several markets already use
(Sec.~\ref{subsec:RelatedWork}).
\dt[confirm the kim2023alscn author list before submission (Nature
Nanotechnology 2023, AlScN/2D-channel BEOL FeFETs)]{}

\subsection{Preliminaries and related work}
\label{subsec:Probes}
\label{subsec:RelatedWork}
\textbf{What a probe reads, and what it costs:} A token is carried
through a transformer by a running vector, the \emph{residual stream},
that every layer reads from and adds back to. After layer $\ell$,
the result is $h^{(\ell)}$: $d$ numbers ($d=8{,}192$ for a 70B
model), one vector per layer, per token. The individual entries are
not labeled features -- no single coordinate means ``deception'' --
so a detector has to read the vector as a whole rather than watch any
one entry. A \emph{linear probe} is the simplest detector over
that object: one stored vector $w$, fitted offline on labeled
activations, and a threshold $\theta$, with
%
\begin{equation}
w \cdot h \;=\; \sum_{i=1}^{d} w_i h_i \;>\; \theta
\;\;\Longrightarrow\;\; \text{flag.}
\label{eq:probe}
\end{equation}
%
Probes of this form recover truth-related structure across several
datasets and model families~\cite{marks2023geometry,burns2023ccs,
zou2023repe}, and the directions are also actionable: adding a multiple
of one back into the residual stream changes behavior, which is usually
called steering~\cite{zou2023repe}. Scoring one activation is $2Ld$ multiply--accumulate
operations for a model of depth $L$ and residual width $d$ -- about
1.3 million operations for a 70B model, or roughly $10^{-5}$ of the
forward pass. This cost is analogous to $L\!\cdot\!d$ rather than
parameter count, and $L\!\cdot\!d$ has not changed significantly
even with frontier models (e.g., a 671B-parameter mixture-of-experts
model has $d=7{,}168$ and $L=61$~\cite{deepseekv3} -- an
$L\!\cdot\!d$ \emph{smaller} than a dense 70B model's
$8{,}192\times80$ -- so a bigger model could have a smaller probe
vector).

\textbf{How detection is graded:} Quality is typically reported as area
under
the receiver operating characteristic curve (\auroc), i.e., the
fraction of (deceptive, honest) pairs the detector orders correctly;
1.0 is perfect and 0.5 is random. An important ``stress test''
for library size is \emph{cross-domain} -- i.e., fit on one kind of
dishonesty and evaluate on another~\cite{ying2026spectrum}.\footnote{A
\emph{domain} is a subject matter to be dishonest about (what a word
means, a checkable fact, a deduction); a \emph{behavior} is a kind of
dishonesty (lying about a fact, flattering the user, quietly
underperforming on an evaluation).}

\textbf{Activation monitoring in production:} The two-stage cascade
in~\cite{ccpp2026,anthropic2026ccblog} and the deployed Gemini
probes~\cite{gemini2026probes} are the closest published statements of
the workload this paper aims to accelerate; both operate at $K$ of
order one.
Monitoring has gotten cheaper as providers moved from reading text to
reading activations: the first-generation text-classifier pair
reported 23.7\% overhead, while the probe-fronted cascade that
replaced it reports 1--3.5\%~\cite{cc2025,ccpp2026}. Separately
reported monitor overheads of roughly 0.1\% of serving
cost~\cite{anthropic2025cheapmonitors} are consistent with watching
just one or two things. What the published figures do not tell us is
how the 1--3.5\% splits between the probe and the escalation stages;
the $2Ld$ arithmetic above bounds the probe's share from below, at
roughly $10^{-5}$ of forward-pass compute.
\dt[still to verify: that the 1--3.5\% in ccpp2026 is measured over
the full cascade rather than the probe alone]{}

\textbf{Approximate nearest-neighbor (ANN) search:}
Locality-sensitive
hashing~\cite{indyk1998lsh,charikar2002simhash} and graph
indexes~\cite{malkov2020hnsw} visit a small fraction of a bank and
still retrieve most true neighbors at sub-millisecond query
times~\cite{faisscuvs2025}. Said methods are likely undesirable
here, as a monitor cannot afford to miss -- an adversary can aim at
the approximation error, while an exhaustive scan has none to
exploit (red-team appendix, App.~A of the companion
document).

\textbf{Associative-search silicon:} Commercially, the nearest analog is
an SRAM-based associative processing unit performing bit-serial Hamming
search at billion-vector scale~\cite{gsi2025apu}. It is an existence
proof for the primitive and the baseline a multi-bit or NVM-based part
would have to differentiate against on density, energy, and metric.
\cam-based accelerators have also been proposed for one-shot
learning~\cite{ni2019fetcam}, nearest-neighbor
search~\cite{kazemi2021fefetcam}, and analog range
matching~\cite{li2020acam}. As such, the hardware requirements of
this workload are not ``new'': the contribution here is the workload
mapping and the library-size analysis rather than a new cell.
\eo[if we keep this paragraph, is there a group SUPREME/JUMP CAM
application-survey citation you want here instead of listing three
individual papers? The deck used the survey figure]{}
% ===== end 02 - Background.tex =====

%%%%%%%%%%%%%%%%%%%%%%%%%%%%%%%%%%%%%%%%%%%%%%%%%%%%%%%%%%%%%%%%%%%%%%
% ===== 03 - Library Scaling.tex =====
%%%%%%%%%%%%%%%%%%%%%%%%%%%%%%%%%%%%%%%%%%%%%%%%%%%%%%%%%%%%%%%%%%%%%%
\subsection{The cost model}
\label{sec:Scaling}
\label{subsec:CostModel}
\label{subsec:Ladder}
Let $K$ be the number of stored signatures per monitored layer, $d$ the
residual width, $L_{\mathrm{mon}}$ the number of monitored layers, $R$
the token rate, and $b$ the bytes per stored element. On a GPU the bank
is read once per monitored token, so
%
\begin{equation}
T_{\mathrm{GPU}} \;=\; K \, d \, b \, L_{\mathrm{mon}} \, R
\quad\text{bytes/s,}
\label{eq:gputraffic}
\end{equation}
%
while a \cam broadcasts the query and leaves the bank in place, so
%
\begin{equation}
T_{\mathrm{CAM}} \;=\; d \, b \, L_{\mathrm{mon}} \, R
\quad\text{bytes/s,}
\label{eq:camtraffic}
\end{equation}
%
independent of $K$. The right-hand axis of Fig.~\ref{fig:ladder}(a)
reads Eq.~(\ref{eq:gputraffic}) at $d=8{,}192$, fp16, one monitored
layer and $R=1{,}000$~tokens/s, against vendor HBM and LPDDR
bandwidths. Monitoring traffic reaches an H100's entire
\qty{3.35}{\tera\byte\per\second}~\cite{nvidia_h100} at
$K\!\approx\!204$K and an A100's at $K\!\approx\!122$K; on a
Jetson-class part at \qty{273}{\giga\byte\per\second} the same
crossing arrives at $K\!\approx\!17$K, roughly an order of magnitude
earlier. Those crossings move with the token rate, which is a
deployment parameter rather than a property of the monitor, so
Fig.~\ref{fig:ladder}(b) plots the limiting $K$ across token rates
spanning batch-1 edge decoding and a loaded datacenter node.

The anchor strip under Fig.~\ref{fig:ladder}(b) places familiar
workloads on that axis. At the per-stream end, a single voice or chat
user decodes at tens of tokens/s, and a deep-reasoning or agentic
stream at a similar instantaneous rate -- what distinguishes those
tasks is duration, at $10^{3}$--$10^{5}$ generated tokens per query
and up to $10^{8}$ tokens per agent run~\cite{aisi2026cyber}.
Per-stream rates batch up: published serving benchmarks for 70B-class
models report aggregate throughputs of order $10^{3}$ tokens/s per
datacenter GPU at high batch~\cite{lyceum2026tps}. At the edge the
two coincide, since a robot-control (vision--language--action, VLA)
policy runs at batch~1 and its stream rate \emph{is} the part's
aggregate rate: an autoregressive VLA at 3--6~Hz and $\sim$7 action
tokens per step is tens of tokens/s, while a chunked $\pi_0$-class
policy at the 19~Hz reported on a Jetson Thor emits 50-action chunks,
i.e., of order $10^{3}$ action tokens/s~\cite{vlaperf2026}, and
continuous video understanding pushes a backbone toward
$10^{3}$--$10^{4}$ token positions/s. As such, the 10--$10^{4}$
tokens/s span of Fig.~\ref{fig:ladder}(b) covers deployed practice
rather than a hypothetical, and the workloads that reach the earliest
crossings -- batch-1 VLA and video on edge parts -- also have no
batch to amortize against.
\dt[anchor numbers to verify before submission: the VLA rows are
converted from VLA-Perf's Hz figures (solid); voice, agentic-session,
and video-backbone rates are order-of-magnitude from usage norms and
frame-rate arithmetic, not from a single citable benchmark. The
serving-node figure still cites a gray-literature roundup -- swap in
the interconnect number if you have one]{}

Two properties of Eq.~(\ref{eq:gputraffic}) justify the work in this
paper. {\bf (1)} Latency is the wrong way to judge a monitor's cost.
At small $K$ the monitor's reads hide under inference latency, so it
\emph{looks} free -- but the bytes and the energy still must move.
Rather, it is better to consider traffic and energy, and
both are charged \emph{per token}. {\bf (2)} A per-token price
multiplies a token base that is compounding: inference volume has
been growing several-fold per year, so a fixed percentage of it is
not a fixed cost. Long-running agents reinforce this
point.
The campaign runs of Sec.~\ref{sec:Introduction} walk 32 steps over
roughly 20 hours at run budgets up to $10^{8}$
tokens~\cite{aisi2026cyber}; a monitor for such a run is always-on by
construction, and its per-token price is what decides whether it
stays on.

% ===== end 03 - Library Scaling.tex =====

%%%%%%%%%%%%%%%%%%%%%%%%%%%%%%%%%%%%%%%%%%%%%%%%%%%%%%%%%%%%%%%%%%%%%%
% ===== 04 - Detector Mapping.tex =====
%%%%%%%%%%%%%%%%%%%%%%%%%%%%%%%%%%%%%%%%%%%%%%%%%%%%%%%%%%%%%%%%%%%%%%
\section{Mapping the Detector into a CAM}
\label{sec:Mapping}

This section covers ways to map activation signatures onto \cam
arrays. Candidate arrays span emerging devices through nearer-term
CMOS solutions that leverage different encoding schemes -- notably,
traditional \cam designs can also be realized with nonvolatile
memories (RRAM, MRAM, etc.), which could approach the density gains
and, with a mostly static signature bank, benefit from
nonvolatility. Sec.~\ref{subsec:Requirements} first spells out what
any \cam-resident realization requires. The first mapping
(Sec.~\ref{subsec:FirstPipeline}) assumed the best-case device -- a
multi-bit, \fefet-class cell whose matchline computes the Euclidean
distance the full-precision detector actually uses -- and it is
retained as the reference point the other encodings are measured
against. The encodings of
Secs.~\ref{subsec:SignRP}--\ref{subsec:TwoStage} satisfy the
requirements on hardware buildable today, including Hamming-distance
options that are near-term and realizable by a host of NVMs as well
as by plain CMOS.

\subsection{What ``CAM-resident'' requires}
\label{subsec:Requirements}
For the monitor to be one memory structure rather than a memory plus an
accelerator, three properties have to hold. {\bf (1)} The query has to
be producible from the raw activation by something cheap, since whatever
sits in front of the array runs on every token. {\bf (2)} The comparison
the array performs natively has to be the comparison the detector needs,
which for a binary or ternary cell means the decision has to be
expressible as a Hamming distance. {\bf (3)} The decision rule has to
consume what a matchline can report -- a winner, a set of rows under a
threshold, or a count -- rather than a full sorted score
vector.\footnote{Returning all $K$ scores is a readout job a \cam is
not for.}

Property~(1) also has a methodological consequence. If the front end is
trained, it can learn to do the classification itself, leaving the
memory as decoration. Sec.~\ref{subsec:Audit} describes the test we
use to detect that, and what it caught when we ran it on pipelines
that did have trained front ends.

\subsection{The first mapping: learned projection, 3-bit Euclidean}
\label{subsec:FirstPipeline}
A multi-bit, \fefet-class cell broadens what a stored row can be:
each cell holds a level rather than a bit, which could allow for a
richer representation of the data -- or the same representation in
fewer memory cells -- with the matchline computing the Euclidean
distance the detector actually uses (Sec.~\ref{subsec:CAMspace}). A
down-projection mechanism does need to be implemented regardless,
since whatever sits in front of the array runs on every token
(requirement~{\bf (1)}): each raw activation is reduced from $d$
dimensions to a much smaller $d'$ and quantized to a few bits per
dimension.
Concretely, the pipeline took the residual-stream
activation at layer 33 (the tuned layer choice described in
Sec.~\ref{subsec:Setup}), built a bank of mutually distinct deception
directions by iterative nullspace
projection~\cite{ravfogel2020inlp} (train a probe, project its direction
out, retrain), weighted each stored direction with one trained scalar,
compressed the activation with a \emph{learned} down-projection from
$d=8{,}192$ to $d'=128$ dimensions, and quantized query and stored rows
alike to 3 bits with learned-step quantization~\cite{esser2020lsq}, so
that the low-precision rounding sits inside the training loop rather
than being applied afterwards.\footnote{Quantization-aware training at
this precision is standard practice: a shipping on-device foundation
model is trained with 2-bit QAT~\cite{apple2025foundation}.} Matching
was Euclidean winner-take-all, returning the winning row and its
margin. $d'=128$ was chosen to match the row geometry of a multi-bit
\fefet array -- 128 cells sharing a matchline, at 3 bits per cell --
so the stored row is a 384-bit-\emph{equivalent} row (physically, 128
cells), which is where the matched 384-bit operating point used
throughout Sec.~\ref{sec:Results} comes
from.

The audit finding that retired this mapping's trained cousins is
reported in Sec.~\ref{subsec:Audit} (Table~\ref{tab:audit}).

\label{subsec:PathBack}%
The encodings that follow trade this mapping's multi-bit cell for
Hamming-distance solutions, which are near-term and realizable by a
host of NVMs as well as by plain CMOS; a frozen projection matched on
a multi-bit Euclidean matchline remains the natural recombination to
try as the device matures.

\subsection{Encoding 1: sign random projection, ternary (or via other
NVMs)}
\label{subsec:SignRP}
A first option for meeting requirements {\bf (1)}--{\bf (3)} on
hardware buildable today is a fixed random projection: draw
$P \in \mathbb{R}^{m \times d}$ once and freeze it, then encode an
activation as the sign pattern of its projection,
%
\begin{equation}
c_j(h) \;=\; \operatorname{sign}\!\big( (Ph)_j \big), \qquad j = 1 \ldots m,
\label{eq:signrp}
\end{equation}
%
with a dead zone $|(Ph)_j| < \epsilon$ mapped to don't-care, which makes
the code ternary and requires a wildcard-capable cell. This is the
standard sign-random-projection (SimHash)
construction~\cite{charikar2002simhash}; the Hamming distance between
two codes estimates the angle between the underlying
vectors.

From a hardware perspective, the comparison is an XOR plus a
population count, which a plain
\cam array computes natively. A code wider than a practical array row
is split across subarrays and the results aggregated, which is
standard practice and a real design parameter -- the measured
\fefet \cam demonstration of~\cite{kazemi2022hdcfefet} uses exactly
this organization, with voting to aggregate results between
subarrays.

\subsection{Encoding 2: thermometer coding (plain binary CMOS)}
\label{subsec:Thermometer}
A second option for reaching a \cam realization -- and one a plain
CMOS binary array can serve: quantize each projected dimension into
$\ell$ levels and encode each
level in unary, so a value of level $v$ becomes $v$ ones followed by
$\ell-v$ zeros. For two codes $u,v$ over the same dimension,
%
\begin{equation}
d_{\mathrm{Ham}}\big(\mathrm{therm}(u), \mathrm{therm}(v)\big)
\;=\; \sum_i \big| u_i - v_i \big|,
\label{eq:thermo}
\end{equation}
%
i.e., Hamming distance on the code is exactly the $L_1$ distance on the
quantized values. The identity is exact, not approximate, which is what
makes the encoding useful here: a cell that can only report disagreement
now supports a magnitude-sensitive metric. No wildcard is needed, so
the cell is a plain binary \cam at roughly 6--10 transistors rather than
a \tcam at roughly 16.

The cell menu is wider than CMOS. A 1\fefet-1C array -- a DRAM-like
cell with the \fefet as access transistor -- has been extended to
ternary \cam operation, with a fabricated proof of concept whose
sensed matchline voltage after charge sharing is linear in Hamming
distance~\cite{fefet1c2026}. A plausible extension is to host either
of this section's encodings on such arrays: the Hamming linearity
reads off the distances of Eqs.~(\ref{eq:signrp})
and~(\ref{eq:thermo}) directly, the don't-care support serves the
ternary code's dead zone, and the DRAM-like cell promises density --
especially considering rising 3D DRAM -- with nonvolatile storage
suited to a mostly static signature bank.
\dt[DAC-thermosensitive.pdf: cite here with its plausible extension
once the PDF arrives from MN's work machine; also settle the
fefet1c2026 bib entry's authors before submission]{}

The comparison against \signrp is therefore not symmetric. At a matched
code length the two encodings use the same number of \emph{bits} but not
the same amount of \emph{silicon}, and Sec.~\ref{sec:Results} reports
both alongside that qualification.

\subsection{The decision rule}

\label{subsec:Decision}
The stored rows are labeled example activations rather than fitted
directions -- i.e., each domain stores not one direction per class
but a few thousand encoded activations, each labeled honest or
deceptive. A query is scored by the fraction of its $k$ nearest stored
rows, under the array's metric, that carry the deceptive label, with
$k=31$ in the runs reported here (the gain flattens well before 31);
\auroc is computed on that score directly, so no threshold is chosen. The scheme is analogous to few-shot classification:
rather than one fitted direction per behavior (the $N$-way, 1-shot
analog), the bank operates $N$-way, many-shot, with thousands of
stored exemplars per domain and the array's distance metric supplying
the comparison. A vote over neighbors rather than a single winner is
worth 3--5 accuracy points in in-distribution tests, and it consumes
only what a best-match or threshold matchline can report.

\subsection{Encoding 3, as a ceiling: shortlist plus exact rerank}
\label{subsec:TwoStage}
To bound what the compression costs rather than only measuring it, we
also run a two-stage variant: stage 1 is a thermometer-Hamming search of
the whole bank that keeps the top $k$ rows as a shortlist, and stage 2
is an exact cosine rerank on the raw, uncompressed activations
restricted to that shortlist. With $k$ equal to the whole bank, stage 2
alone defines a full-precision ceiling for the same data and protocol.
Shortlist-then-rerank is standard practice in large-scale retrieval
(e.g., the two-stage GPU indexes of~\cite{faisscuvs2025});
equivalently, stage 2 can be read as checking a smaller $K$ that a
GPU can easily and accurately serve, with stage 1 reducing the
library to that size. The variant is also the architecture a real
system might adopt if
single-stage fidelity is inadequate, at the cost of a second memory
structure and a floating-point stage.
% ===== end 04 - Detector Mapping.tex =====

%%%%%%%%%%%%%%%%%%%%%%%%%%%%%%%%%%%%%%%%%%%%%%%%%%%%%%%%%%%%%%%%%%%%%%
% ===== 05 - Results.tex =====
%%%%%%%%%%%%%%%%%%%%%%%%%%%%%%%%%%%%%%%%%%%%%%%%%%%%%%%%%%%%%%%%%%%%%%
\section{Workload 1: Deception Detection}
\label{sec:Results}

We first evaluate the mappings of Sec.~\ref{sec:Mapping} on a
deception detection workload, where the
library is set by how many behaviors and domains someone decides to
watch -- here, the $K\approx10^{2}$--$10^{3}$ regime of
Fig.~\ref{fig:ladder}. We describe the experimental setup, then
report {\bf (1)} the ten-domain fidelity of each encoding, {\bf (2)}
whether wider codes improve cross-domain transfer, {\bf (3)} how the
two-stage shortlist behaves across bank sizes, and {\bf (4)} the
destroyed-vector audit that disciplines every configuration reported.
Sec.~\ref{sec:LargeK} evaluates the other end of the $K$ axis, where
the library is set by the vocabulary.

\begin{figure*}[t]
\centering
\includegraphics[width=7.0in]{fig_grids_v3.pdf}
\caption{Cross-domain \auroc on the ten-domain grid,
Llama-3.3-70B-Instruct activations at layer 33. Rows index the domain
whose examples are stored in the bank, columns the domain the held-out
test activations come from. (a) 384-bit \signrp ternary code, which
requires a wildcard-capable \tcam cell; (b) 384-bit thermometer code,
which requires only a plain binary \cam cell; (c) Hamming shortlist
over the whole bank followed by an exact cosine rerank on the raw
activations, i.e., the full-precision ceiling for the same data and
protocol.}
\label{fig:grids}
\end{figure*}

\subsection{Experimental setup}
\label{subsec:Setup}
Our evaluation case follows the cross-domain protocol of the
Truthfulness Spectrum study~\cite{ying2026spectrum}: a model produces
honest and deceptive statements across several subject-matter
domains, its activations are recorded, and a detector built from one
domain's examples is scored on every other domain's -- the ``stress
test'' of Sec.~\ref{subsec:Probes}. Activations are taken from
Llama-3.3-70B-Instruct~\cite{llama3herd}. The evaluation grid has ten
domains:
definitional, empirical, fictional, logical, ethical, sycophancy,
\texttt{repe\_honesty}, roleplaying, insider trading, and sandbagging. The
first four plus sycophancy correspond to the published
grid~\cite{ying2026spectrum}; \texttt{repe\_honesty} follows the
representation-engineering dataset~\cite{zou2023repe}. Within each
domain, 70\% of examples form the stored bank and 30\% are held out for
test. Every cell of the grid is a cross-domain measurement: rows index
the domain whose examples are in the bank, columns the domain the test
activations come from. Bank sizes differ substantially by domain, from
831 rows (definitional) to 3{,}201 (sycophancy), which turns out to
matter in Sec.~\ref{subsec:Shortlist}. Activations are read at layer
33: the source study selected that layer (with logistic-regression
probes on average-token activations) from a cross-domain tuning sweep
on its truth-type datasets, and we inherit the choice so our grids
can be compared against theirs -- i.e., it is a tuned choice for this
model family and scale rather than a transferable constant. All
numbers reported below are from cluster runs on this
configuration.

\subsection{The grids}
\label{subsec:Grids}

Fig.~\ref{fig:grids} shows the ten-domain cross-domain grids
for the two single-stage encodings and for the exact-rerank ceiling;
Table~\ref{tab:encodings} summarizes them alongside the hardware each
requires. Each grid is read the same way: a row names the domain
whose examples are stored in the bank, a column the domain the
held-out test activations come from -- i.e., a cell scores a bank
built on behavior/domain $X$ against test material from
$Y$.

We make three observations on the ten-domain grids.

\textbf{In-domain fidelity is essentially unaffected by the encoding:}
When a bank is tested on the domain it was built from, the compressed
codes give up almost nothing: the diagonal means are 0.920 (\signrp
ternary, 384~b) and 0.919 (thermometer, 384~b), against 0.928 for the
exact rerank. Said differently, replacing full-precision cosine
similarity over the raw 8{,}192 values with a 384-bit code in a plain
array costs 0.008--0.009 \auroc. The factual block (definitional,
empirical, fictional -- all nine cells, transfer included) averages
0.974 for ternary and 0.942 for thermometer, against the 0.978
ceiling.

\textbf{Cross-domain fidelity is where compression shows, and it
separates the two encodings:} An off-diagonal cell scores a bank
built from one domain against test examples from another, so it
measures how well a stored library transfers to material it was not
built from. That transfer is where the encodings pay: off-diagonal
means are 0.640 for ternary and 0.579 for thermometer, against 0.705
for the full-precision ceiling -- a cost of 0.065--0.126 \auroc,
where in-domain the same compression cost less than 0.01. The
ordering is the surprise. At this matched 384-bit point, ternary is
ahead of thermometer by 0.061 \auroc on average and in 63 of 90
cells, which is the opposite of what a preference for the cheaper
cell would predict. One caveat keeps the question open: the
comparison is matched in \emph{bits} but not in \emph{silicon},
since a ternary cell needs wildcard capability. If a ternary position
really costs two bits of hardware, a matched-silicon comparison would
give thermometer roughly twice the code length; that experiment has
not been run.
\dt[R2, per IB: fold in the matched-cell-count run -- thermometer at 768~b scores 0.924 in-domain / 0.631 cross-domain against ternary-384's 0.920 / 0.640, ternary ahead in 54 of 90 cells (was 63), worst cell 0.319 vs 0.259 -- so at matched cell count the two are within noise on the mean; replace the open-question wording above]{}

\textbf{The pattern of failure reproduces the published one:} All three
grids keep the factual block high and put sycophancy near chance from
every other row, which is the qualitative result the source study
reports~\cite{ying2026spectrum}. Sycophancy's own in-domain cell is the
weakest diagonal entry in every configuration (0.627, 0.633, 0.636),
stable across conditions, so it appears to be a property of the data
rather than a compression artifact. The weakest off-diagonal cell sits
at a similar depth in every configuration (0.259, 0.226, 0.220) but
its \emph{location} moves with the encoding:
\texttt{repe\_honesty}~$\to$~sandbagging for ternary,
logical~$\to$~insider trading for thermometer, and
\texttt{repe\_honesty}~$\to$~fictional for the rerank
configurations. The bigger picture is that the accuracy ceiling on
transfer is set by the representations, not by our compression -- the
full-precision grid has the same holes, in the same places. Covering
a behavior suite at this accuracy therefore means storing directions
for each behavior and domain rather than hoping one direction
travels, i.e., operating at the library scale this paper
prices.

\begin{table}[t]
\centering
\caption{The encodings, their measured fidelity on the ten-domain grid,
and what each needs from a cell. Bank size is for a 4{,}000-row exemplar
bank at the given code width. The in-domain minimum is sycophancy in
every row (0.627 / 0.633 / 0.640 / 0.636 top to bottom); the
cross-domain minimum is 0.259 / 0.226 / 0.219 / 0.220 top to bottom,
and its cell location varies by encoding
(Sec.~\ref{subsec:Grids}).}
\label{tab:encodings}
{\scriptsize\setlength{\tabcolsep}{3.5pt}
\begin{tabular}{@{}lccll@{}}
\toprule
\textbf{Encoding} & \textbf{in-dom.} & \textbf{cross-dom.} &
\textbf{Cell needed} & \textbf{Bank} \\
 & mean & mean & & (4K rows) \\
\midrule
\signrp ternary, 384\,b       & 0.920 & 0.640 & \tcam, $\sim$16T   & 192\,kB \\
Thermometer, 384\,b           & 0.919 & 0.579 & BCAM, $\sim$6--10T & 192\,kB \\
Shortlist $+$ rerank          & 0.928 & 0.671 & BCAM $+$ FP unit   & 192\,kB $+$ raw \\
Exact rerank (ceiling)        & 0.928 & 0.705 & FP baseline        & 65\,MB \\
\bottomrule
\end{tabular}}
\end{table}

\subsection{Cross-domain transfer is hard, and wider codes do not
close the gap}
\label{subsec:BitWidth}

The transfer gap of Sec.~\ref{subsec:Grids} raises the natural
question of whether more code capacity would close it -- and the
answer matters, because if it does not, covering more behaviors means
storing \emph{more} comparison vectors rather than wider ones, i.e.,
growing $K$. Widening the code should tighten the angle estimate in
Eq.~(\ref{eq:signrp}) and improve matching, and in-domain it does.
Cross-domain it does not, in either encoding. Growing the code from 384
to roughly 2{,}048 bits, 23 of 90 off-diagonal cells drop by more than
0.05 \auroc for the ternary code, with a worst single-cell change of
$-0.303$ (sandbagging tested on insider trading); for the thermometer
code the same growth moves 12 of 90 cells by more than 0.05, worst
$-0.209$. Cell by cell, for the ternary code, the diagonal is flat or
slightly up, the losses concentrate in the rows that transfer worst at 384 bits
(\texttt{repe\_honesty}, sandbagging, sycophancy), and the one column
that improves broadly -- insider trading -- is a domain, not the
encoding, getting easier to hit.

A plausible reading is that extra dimensions preserve more of the
geometry of the training domain, including the parts that do not
generalize, so a wider code is a more faithful encoding of a
domain-specific decision boundary. Whatever the mechanism, the
takeaway for hardware is simple: {\bf (1)} the cheaper array is also
the one that transfers better -- on this grid, the costlier cell does
not buy accuracy -- so there is no argument for
paying more silicon per row; and {\bf (2)} if transfer must come from
somewhere, it comes from storing more vectors, not wider ones --
which is the library growth Sec.~\ref{sec:Scaling} prices.

\subsection{A fixed shortlist does not travel across bank sizes}
\label{subsec:Shortlist}
The two-stage variant of Sec.~\ref{subsec:TwoStage} is the natural way
to recover the ceiling, and at first reading it works: with an absolute
shortlist of $k=400$, the diagonal mean is unchanged at 0.928 and the
off-diagonal mean falls only from 0.705 to 0.671. The aggregate hides a
failure of a \emph{fixed} parameter, which matters more for a hardware
design than for a software one.

Bank sizes differ by roughly $4\times$ across domains, from 831 rows
(definitional) to 3{,}201 (sycophancy). Coverage -- how often the true
full-precision top-1 neighbor actually lands inside the Hamming
shortlist -- shows the consequence. For the definitional bank
the mean is 82.5\%, against a chance floor of $k/N = 48.1\%$. For the
sycophancy bank the mean is 18.3\% against a floor of 12.5\%, i.e.,
barely above chance, and several test domains sit at or below their own
floor. Consequently, 8 of the 10 training rows are stable between $k=400$
and full precision while the sycophancy row loses 0.204 \auroc on
average and up to 0.468 in a single cell; the ethical row shows a smaller
version of the same effect at $-0.079$.

This is a sizing problem rather than a broken mechanism. Sweeping $k$ as a fraction of each domain's own bank,
sycophancy's cross-domain mean recovers from 0.532 at 12.5\% to 0.736 at
100\%, mostly monotonically, and the 100\% value matches the
independently computed full-precision row average to three decimal places
(0.736 vs.\ 0.736; definitional 0.750 vs.\ 0.751), which is a useful
self-consistency check on the pipeline. The recovery costs 75--100\% of
the bank, at which point it is not a shortlist. For a \cam-resident
design, that argues for the single-stage encodings above, or for a
threshold rather than a top-$k$ criterion so the number of survivors
adapts to the bank instead of being fixed at design time.

\subsection{The destroyed-vector audit}
\label{subsec:Audit}
\label{subsec:AuditResults}
\begin{table}[t]
\centering
\caption{The destroyed-vector audit across five variants of the first
pipeline (July 30). Front-end share is the accuracy still present after
every stored row is replaced by a random vector, i.e., the share the
trained layers were carrying alone. High is bad: 100\% means the memory
contributed nothing. This audit is our sanity check that a pipeline is
not ``cheating,'' i.e., classifying in a trained front end rather than
in the stored bank; we run it on every configuration reported in this
paper.}
\label{tab:audit}
{\footnotesize
\begin{tabular}{@{}lcl@{}}
\toprule
\textbf{Variant} & \textbf{Front-end share} & \textbf{Outcome} \\
\midrule
Learned binary hash          & 98--100\% & retired \\
Weighted thermometer         & 66--85\%  & retired \\
3-bit LSQ, Euclidean         & 24--28\%  & survives \\
Direct sign, no projection   & 20\%      & survives \\
Random sign projection       & 4--14\%   & survives \\
\bottomrule
\end{tabular}}
\end{table}
Any pipeline of the form ``trained front end, then memory'' can pass
its accuracy test while the memory contributes nothing. We therefore
re-measure every configuration with the stored rows replaced by random
vectors of the same shape: a \cam-based detector should collapse to
chance when the bank is destroyed, and if it does not, whatever sits
in front of the array was doing the classification. We report the
accuracy that survives destruction as the \emph{front-end share},
where high is bad: 100\% means the memory was decoration.

The audit has changed decisions. Table~\ref{tab:audit} reports it
across five variants of the first pipeline
(Sec.~\ref{subsec:FirstPipeline}). In the two highest-scoring
variants, one of which reached 0.993 \auroc, between 66\% and 100\%
of the accuracy survived randomizing every stored row: the trained
layers in front of the array were performing most of the
classification, and the \cam was close to decoration. Both were
retired. The 3-bit LSQ Euclidean configuration sits at 24--28\%, so
its memory is contributing; the variants whose stored rows carry
nearly everything are the ones with a frozen or absent projection, at
4--20\%, which is why the encodings of Secs.~\ref{subsec:SignRP}
and~\ref{subsec:Thermometer} use frozen projections -- with nothing
trained in front of the array, the failure has nowhere to occur.

The encodings reported above pass. Randomizing every stored
activation collapses the \signrp grid to a mean of 0.50, with 83\%
of cells between 0.4 and 0.6, so the classification is coming from
what is stored rather than from anything in front of the array. An
accuracy number alone would not have separated the five variants of
Table~\ref{tab:audit} -- the highest-scoring one was the worst
offender. As such, we now run the audit by default on every
configuration, report front-end share alongside \auroc, and treat a
variant whose accuracy survives bank destruction as an accelerator
result rather than a memory result. The recombined design of
Sec.~\ref{subsec:PathBack} should pass this audit by construction --
nothing is trained ahead of its array -- and will be checked the
same way when it is run.

% ===== end 05 - Results.tex =====

%%%%%%%%%%%%%%%%%%%%%%%%%%%%%%%%%%%%%%%%%%%%%%%%%%%%%%%%%%%%%%%%%%%%%%
% ===== 06 - Large K.tex =====
%%%%%%%%%%%%%%%%%%%%%%%%%%%%%%%%%%%%%%%%%%%%%%%%%%%%%%%%%%%%%%%%%%%%%%
\section{Workload 2: Whole-Vocabulary Readout}
\label{sec:LargeK}

This section evaluates the mappings on the second workload:
whole-vocabulary readout, where the library is set by the model's
vocabulary rather than by anyone's choices -- the
$K\approx1.3\times10^{5}$ rung of Fig.~\ref{fig:ladder}. We describe
the published method and why its $K$ is structural, report an
in-house reproduction at watchlist scale with its measured fidelity
at \cam precision, and price what full readout would cost a
\cam-resident design against streaming the same bank from HBM.

\subsection{A workload whose $K$ nobody selected}
\label{subsec:Jspace}
The grids in Sec.~\ref{sec:Results} are representative of the
$K\approx10^{2}$--$10^{3}$ regime, where the library grows with the
number of behaviors somebody thinks to watch. A recently published
readout method sets $K$ a different way. It reports that a model holds words ``in
mind'' without emitting them, and that for every token in the vocabulary
there is an internal pattern meaning that word is currently active,
readable whether or not the word is ever
said~\cite{anthropic2026workspace}. In the spider example of
Sec.~\ref{sec:Introduction}, the word appears in neither prompt nor
answer, and replacing the internal \emph{spider} pattern with the one
for \emph{ant} changes the answer from eight to six. The same work
reports that the readout catches misbehavior beyond the staged safety
test: ``manipulation'' was active while a model produced falsified
data during a pre-release audit.

For our purposes the relevant property is structural rather than
behavioral. Reading this space means comparing the internal state against
one stored pattern \emph{per vocabulary token} -- roughly 128K signatures
for a Llama-family tokenizer -- at several layers, every token. Nobody
chose that library size; it falls out of the vocabulary. And it is
plausibly a floor, since the published method tracks single words, and
phrases would only enlarge the library.

\subsection{An in-house reproduction, at watchlist scale}
\label{subsec:Watchlist}
We reproduced the readout in-house on Llama-3.1-8B: a layer-by-layer
ranked list of what the model is currently representing. About 13 of 32
layers carry a usable readout at this scale, which is the source of the
``13--25 usable layers'' factor used in Sec.~\ref{sec:Scaling}. Each
signature is precomputed offline, so a watched word costs one comparison
per token, and we stored 47 deception-adjacent signatures per monitored
layer as \cam rows.

In a staged scenario in which the model is told it is about to be
replaced and is given a personnel file it was not meant to use, ten watch
words were tracked by exact full-vocabulary rank across the 2{,}048
prompt tokens the model reads before generating anything. Five of them --
``scandal,'' ``evidence,'' ``leaked,'' ``fake'' and ``shutdown'' --
reached rank~0 at interpretable moments in the reading, ``threat''
reached rank~1 and ``blackmail'' rank~2, and none of the ten appears in
the prompt text. The reply the model then produced was cooperative and
unremarkable, i.e., not something a text classifier would score as
harmful. Ground truth for the reproduction is the model's own
full-precision readout, so no labeled deception data is required
anywhere.

These are preliminary numbers: one model, one staged scenario, ten watch
words in the demonstration and 47 stored signatures.\footnote{Several of
the fidelity and scaling values in this subsection were read from result
charts rather than from result files; verification against the team's
files is queued, and the scaling sweep re-runs with three known fixes
before any of it is reused externally.}

\subsection{Fidelity at \cam precision}
\label{subsec:JspaceFidelity}
The deception grid of Sec.~\ref{sec:Results} needs labeled activations.
Readout does not: the full-precision, full-dimension dot product against
the stored patterns \emph{is} the ground truth, so a compressed
implementation can be scored directly against the exact readout it is
trying to reproduce. Compression can therefore be measured here without
labels, and the deception numbers of Sec.~\ref{sec:Results} are not
assumed to carry over to it.

The setup takes J-space vectors at a fixed layer, gathers residual-stream
activations from Wikitext, records which stored vectors fire above a
global threshold under the exact computation, and asks how well each
\cam-compatible encoding reproduces that list. Two readouts are scored:
\emph{threshold}, where every row that fires is reported and the metric
is per-row firing F1, and \emph{top-1}, where the single closest row is
reported and the metric is agreement with the exact readout's choice.
The bank is 47 deception-adjacent vectors per layer, as in
Sec.~\ref{subsec:Watchlist}.

Fig.~\ref{fig:jspace}(a,b) plots both against storage bits per \cam
row, which is the axis that connects fidelity to
Sec.~\ref{sec:Hardware}: bits per row set the bank size, and the bank
size sets what has to be held or moved. An uncompressed fp16 row at
$d=8{,}192$ is 131{,}072 bits and reproduces the exact readout by
definition. Three observations.

\textbf{Two orders of magnitude of compression cost less than it might:}
A 3-bit Euclidean code at $\sim$$10^{3}$ bits per row -- a $128\times$
reduction -- reaches \hbox{0.96} top-1 agreement and 0.81 per-row firing
F1. Keeping the projection in full precision and quantizing nothing but
the width (the exact-$L_2$ series) reaches 1.00 top-1 at 2{,}048 bits and
0.84 firing F1 at 4{,}096, so most of the remaining gap at $10^{3}$ bits
is the quantizer rather than the projection.

\textbf{The top-1 readout is markedly cheaper than the threshold
readout:} At every code width the top-1 curve sits above the firing-F1
curve, by 0.15--0.2 in the range that matters. Reporting \emph{which row
is closest} survives compression better than reporting \emph{which rows
cross a threshold}, because the first is a comparison among rows and the
second is a comparison against an absolute value that quantization
shifts. For a \cam this is a matchline-mode result: best-match mode
tolerates a cheaper cell than threshold mode does, at equal fidelity.

\textbf{Hamming and thermometer codes do not reach the same operating
point here:} Both rise with code width to roughly 0.74 top-1 and 0.49
firing F1 at $10^{3}$ bits, and the Hamming series then falls sharply at
4{,}096 bits, where the encoding switches to a direct sign hash without a
learned projection. That is the opposite ordering from
Sec.~\ref{sec:Results}, where at a matched 384 bits the sign code was
ahead. The two workloads are asking different questions -- separating two
labeled classes there, reproducing a specific ranked list here -- and the
metric a readout needs appears to be the one a multi-bit cell provides.
Sec.~\ref{subsec:Devices} takes up what this would mean for the
array.
\dt[R2, per IB: the recombination run answers the projection-artifact
question -- frozen RP to 128, 3-bit, Euclidean matchline 0.920 / 0.579
vs Hamming 0.920 / 0.593 (unrounded Euclidean 0.916 / 0.598), so once
the projection is frozen the multi-bit cell buys no accuracy, only
fewer cells per row; fold in here and in VII-D]{}

\dt[Fig.~5(c) is orphaned by the VI-D cut: regenerate fig\_jspace
without panel (c) and at reduced height, or restore two sentences on
the library-size sweep]{}
\begin{figure*}[t]
\centering
\includegraphics[width=7.0in]{fig_jspace_v1.pdf}
\caption{Whole-vocabulary readout reproduced through a \cam.
(a) Per-row firing F1 and (b) top-1 agreement with the exact J-lens,
against storage bits per \cam row, for a bank of 47 vectors per layer;
the star marks an uncompressed fp16 row (131{,}072 bits), which
reproduces the exact readout by definition. (c) Per-row macro F1 as the
library grows from $10^{2}$ to $10^{4}$ tracked vectors at 124$+$4
cells per row. Values digitized from the source deck.}
\label{fig:jspace}
\end{figure*}

\subsection{What full readout would cost}
\label{subsec:ReadoutCost}
Priced at 70B geometry, an uncompressed bank of 128K signatures at
$d=8{,}192$, fp16, is \qty{2.1}{\giga\byte} per layer. Streaming that on
every token at 1{,}000 tokens/s over 25 monitored layers is roughly
\qty{53}{\tera\byte\per\second}, which is not a workload anyone runs. The
useful part of the answer, however, stays small: only on the order of 25
patterns are meaningfully active at any one time, so the operation is
detection rather than readout -- scan everything, return the few that
fired -- which is the question a threshold-mode \cam answers natively.

Three things would make it tractable, in increasing order of what they
require.

\textbf{A watchlist:} Store the few hundred words already known to be
worth watching, which drops $K$ by three orders of magnitude. Our 8B
demonstration in Sec.~\ref{subsec:Watchlist} is this case. A watchlist
catches only what somebody thought to list, and the words that matter
most may be the ones nobody anticipated, so it is a first rung toward
full readout rather than a substitute for it.

\textbf{Threshold matching instead of ranking:} A \cam can report which
rows are within $\tau$ without scoring the rest, and with roughly 25
active patterns out of 128K nearly every query returns nothing.

\textbf{Compressing the stored side:} At the code widths
Sec.~\ref{sec:Results} measured, 384--2{,}048 bits per row against
131{,}072 bits for an uncompressed fp16 row, the whole 128K bank is
6--\qty{33}{\mega\byte} per layer (Fig.~\ref{fig:capacity}), which fits
on a die. The requirement is that it not move. Only the first of the three has been
demonstrated here; the other two are design arguments.

\begin{figure}[t]
\centering
\includegraphics[width=3.4in]{fig_capacity_v1.pdf}
\caption{Bank capacity per monitored layer against library size, for
uncompressed fp16 rows and for the two code widths measured in
Sec.~\ref{sec:Results}, with two on-chip capacity references. A
whole-vocabulary bank at the measured code widths is
6--\qty{33}{\mega\byte} per layer; the same bank uncompressed is
\qty{2.1}{\giga\byte}.}
\label{fig:capacity}
\end{figure}
% ===== end 06 - Large K.tex =====

%%%%%%%%%%%%%%%%%%%%%%%%%%%%%%%%%%%%%%%%%%%%%%%%%%%%%%%%%%%%%%%%%%%%%%
% ===== 07 - Hardware Analysis.tex =====
%%%%%%%%%%%%%%%%%%%%%%%%%%%%%%%%%%%%%%%%%%%%%%%%%%%%%%%%%%%%%%%%%%%%%%
\section{Hardware Implications}
\label{sec:Hardware}

\subsection{Monitoring as a share of the host's own traffic}
\label{subsec:Tax}
Eq.~(\ref{eq:gputraffic}) is easier to interpret against what the model
itself already moves. In the datacenter, a 70B model at fp8 reads roughly
\qty{70}{\giga\byte} of weights per token; a library of $K=100$K rows at
$d=8{,}192$ adds \qty{1.6}{\giga\byte}, or about 2\%. At the edge, an 8B
model at 4-bit reads roughly \qty{4}{\giga\byte} per token, and $K=25$K
at $d=4{,}096$ adds \qty{205}{\mega\byte}, about 5\%, rising to roughly
20\% by $K=100$K. The same architectural point therefore arrives one to
two orders of magnitude earlier in $K$ at the edge, on parts that have no
HBM available at any price. Edge inference is already bandwidth-limited,
so every monitoring byte trades approximately 1:1 against token rate, and
the 15--\qty{60}{\watt} envelope makes DRAM energy per bit a first-order
term rather than a footnote.

\subsection{The cache cliff, and how close current banks are to it}
\label{subsec:CacheCliff}
Capacity produces a sharper discontinuity than bandwidth. Consider a
Jetson-class part: \qty{273}{\giga\byte\per\second} LPDDR5X, roughly
\qty{4}{\mega\byte} of shared L2 and no HBM~\cite{nvidia_thor}, running
an 8B assistant onboard and monitoring only the four behaviors observed
in the campaign of Sec.~\ref{sec:Introduction}. At $d=4{,}096$, fp16, 500
stored rows is \qty{4.1}{\mega\byte}, which is that shared L2 almost
exactly. Below the cliff the bank is resident and monitoring is close to
free; above it the bank streams from LPDDR on every token.

The arithmetic of Sec.~\ref{sec:Introduction} put four behaviors at
520--1{,}000 rows, and the banks currently in hand are already in that
range: 47 signatures across 13 layers is roughly 600 rows, and the
exemplar bank of Sec.~\ref{sec:Results} is about 4{,}000. The fair
pushback is that at the code widths Sec.~\ref{sec:Results} measured,
today's banks compress back under the cliff -- 4{,}000 rows at 2{,}048
bits is about \qty{1}{\mega\byte}. Two qualifications apply. Compression has been validated for detection, not for the
steering-grade full-precision directions that a correcting monitor would
need. And the cliff does not move as the library grows: at
whole-vocabulary scale even the compressed bank is out of cache
(Fig.~\ref{fig:capacity}). The edge crossing is therefore not purely a
$10^{5}$ hypothetical; the libraries already in hand sit near the
threshold.

\subsection{Where the search engine could sit}
\label{subsec:Placement}
Table~\ref{tab:placement} lists the placements and what each trades. The
trade is always the same: how far the query has to travel, and what it
has to contend with on the way. A monitor co-resident on the inference
GPU competes with live inference for cache and bandwidth, which is the
contention the analysis above assumes away and which would need to be
measured. A back-end-of-line (BEOL) search engine in the metal layers
above GPU compute keeps the query travel to microns and the library
nonvolatile and static, so it augments a part already purchased -- but no
foundry PDK offers BEOL \fefet today, and BEOL RRAM/MRAM at 22/16\,nm is
the defensible near-term citation
(the BEOL device work is in Sec.~\ref{subsec:CAMspace}), which makes
that row a 2030s option and it should be labeled that way.
\dt[confirm you want the BEOL row in a DATE paper at all. It is the most
speculative entry in the table and DATE reviewers are less forgiving
than a sponsor audience. The CXL/appliance row is the one an industrial
reviewer can price today]{}

\begin{table}[t]
\centering
\caption{Placements for a \cam-resident monitor.}
\label{tab:placement}
{\footnotesize
\begin{tabular}{@{}p{0.27\columnwidth}p{0.66\columnwidth}@{}}
\toprule
\textbf{Placement} & \textbf{Trade-off} \\
\midrule
Co-resident on the inference GPU &
Contends with live inference for cache and bandwidth; the contention
term is not in our numbers and could dominate at scale. \\
Dedicated monitoring GPU(s) &
Whole GPUs scanning a static table, plus cross-device latency. \\
\cam or CXL-attached appliance &
Matches in place, but the query crosses a board-level bus; a separate
part to build and qualify. \\
BEOL search engine above GPU compute &
Query travels microns, no cache contention, nonvolatile static library;
no PDK support for BEOL \fefet today. \\
\bottomrule
\end{tabular}}
\end{table}

\subsection{Device candidates}
\label{subsec:Devices}
Table~\ref{tab:devices} lists candidate substrates. The encodings of
Sec.~\ref{sec:Mapping} matter here because they change which row is
required. A thermometer code needs only the last row, i.e., no new
device: a binary SRAM \cam in any CMOS process. The \signrp ternary code
needs wildcard capability, available in CMOS \tcam at higher area or in a
2-transistor \fefet cell~\cite{ni2019fetcam}. A multi-bit Euclidean
matchline is the row that requires a multi-level
device~\cite{kazemi2021fefetcam}. The two-stage variant would benefit
from it, as would the first pipeline of Sec.~\ref{subsec:FirstPipeline},
whose 3-bit Euclidean matching maps directly onto such a cell, and the
frozen-projection Euclidean design of Sec.~\ref{subsec:PathBack}. The
multi-bit row therefore stays in the table even though the encodings
this paper measures do not require it: it is the substrate for the next
experiment rather than a device the current results depend on.

\textbf{The two workloads may want different cells:} The measured
results give the multi-bit row a second reason to stay. At a matched
384 bits the \signrp ternary code is ahead of thermometer on the
deception grid (Sec.~\ref{sec:Results}), while on the readout workload
the 3-bit Euclidean code is well ahead of every Hamming-metric encoding
at the same storage budget (Sec.~\ref{subsec:JspaceFidelity}). The two
workloads are asking different questions -- separating two labeled
classes there, reproducing a specific ranked list here -- and the
ordering flip is consistent with the readout needing the magnitude
information that sign codes discard. If the flip survives two
checks -- matched cell count, and a common projection across encodings
-- it is a design-space result: a monitoring engine that has to serve both ends of
Fig.~\ref{fig:ladder} either provisions two cell types or pays for the
multi-bit cell and serves the small-$K$ workload with it too. If it
does not survive them, one array serves both. Write endurance appears in
the table because the bank is rewritten whenever the monitored model is
updated or the probe layer is re-tuned, which is a per-model-release
event rather than a per-token one; that is the property that makes an NVM
bank reasonable here and unreasonable for transformer intermediates.

\begin{table}[t]
\centering
\caption{Candidate substrates for the signature bank. Device figures are
literature values, not our measurements.}
\label{tab:devices}
{\footnotesize
\begin{tabular}{@{}lcccl@{}}
\toprule
\textbf{Tech} & \textbf{Bits/cell} & \textbf{Write} & \textbf{Endurance} &
\textbf{Best fit here} \\
\midrule
\fefet   & 2--3   & $\sim$20\,ns    & $>$5$\times$10$^{10}$ & multi-bit Euclidean \\
ECRAM    & up to 6 & $\sim$200\,$\mu$s & $\sim$10$^{12}$    & static direction bank \\
RRAM     & 1--3   & sub-$\mu$s      & $\sim$10$^{9}$        & density; wear-leveling \\
SRAM \cam & 1     & SRAM            & --                    & Hamming; any CMOS fab \\
\bottomrule
\end{tabular}}
\end{table}

\subsection{A first-order energy and latency estimate}
\label{subsec:CostModelGap}
\dt[R2, per IB: measured H100 PCIe baselines at 131K rows -- fp16 bank
1.22~ms / 277~mJ per token at batch 1, 21~$\mu$s / 4.7~mJ at batch 64
(achieved 1.6--1.8~TB/s, so the 204K crossing is ${\approx}$107K in
practice); FAISS popcount on the 384-b codes 28~$\mu$s / 2.1~mJ at
batch 1, 0.5~$\mu$s / 0.04~mJ at batch 64 -- go into Table~VI and this
paragraph. Framing needs MN: compression alone brings a GPU to
${\approx}$3\% of its per-token budget at $10^{3}$~tok/s, so the CAM case
rests on search energy and non-contention]{}
The analysis above is capacity and bandwidth arithmetic. A full cost
model against a tuned GPU baseline is future work; what can be stated
now, from
published device and memory figures, is a first-order estimate of the
energy and latency of the search itself, and the gap it indicates is
wide enough to survive considerable error in the inputs.

On the GPU side, the dominant energy term at large $K$ is reading the
bank out of DRAM. Measured HBM2 access energy is approximately
3.97~pJ/bit~\cite{oconnor2017fgdram}. At the whole-vocabulary
operating point -- $K=1.3\!\times\!10^{5}$ rows of fp16 at $d=8{,}192$,
i.e., the \qty{2.1}{\giga\byte} bank of
Sec.~\ref{subsec:ReadoutCost} -- reading the bank once is roughly
\emph{67~mJ per token per layer}, before any arithmetic is charged. On
the \cam side, the same search is $K$ rows of 2{,}048-bit code compared
in place, about $2.6\!\times\!10^{8}$ cell comparisons. Published
figures for binary and ternary \cam search sit in the range of
0.1--1~fJ per cell per search across CMOS survey data and fabricated
\fefet arrays~\cite{pagiamtzis2006cam,ni2019fetcam,liu2022evacam},
which puts the whole-bank search at roughly \emph{0.03--0.3~$\mu$J per
token per layer} -- five to six orders of magnitude under the DRAM
read. The gap decomposes into two factors: $64\times$ from the code
width (2{,}048 bits stored per row against 131{,}072; measured
fidelity at that compression is what Secs.~\ref{sec:Results}
and~\ref{sec:LargeK} establish) and $\sim$$4\!\times\!10^{3}$ from the
substrate (fJ-scale cell comparison against pJ-scale DRAM bit
movement). As such, even a $10\times$ error in the per-cell figure --
which spans the published range -- moves the conclusion by one order
out of five or six. Scaled to a deployment, the same arithmetic at 25
monitored layers and $10^{3}$~tokens/s is roughly 1.7~kW of DRAM
access power on the GPU side against tens of milliwatts of array
search: monitoring at this $K$ is not a workload a deployed part
absorbs, which is what Fig.~\ref{fig:ladder}(b) says in bandwidth
terms.\footnote{What the estimate omits, on each side: the GPU figure
excludes the distance arithmetic itself and any cache reuse across
batched requests; the \cam figure
takes array-level published values, which include matchline and
sense-amp energy but not the query projection ($\sim$$10^{6}$
multiply--accumulates per token, negligible at this scale) or the
mismatch-count reduction across subarrays. These are first-order
literature values, not our measurements.}

Latency divides the same way. Streaming \qty{2.1}{\giga\byte} at an
H100's \qty{3.35}{\tera\byte\per\second} is a 0.63~ms floor per token
per layer at batch~1, or roughly 10~$\mu$s per request amortized at
batch 64. Fabricated and modeled TCAM arrays report search times of
0.3--4.4~ns~\cite{liu2022evacam,ni2019fetcam}; with subarrays searched
in parallel and mismatch counts reduced in a tree, a whole-bank search
plausibly completes in tens of nanoseconds, and even at 100~ns it sits
two orders under the amortized floor and roughly four under the
batch-1 floor. Table~\ref{tab:costmodel} carries these entries, marked
as estimates, alongside the rows this paper measures and the rows that
remain open.

\begin{table}[t]
\centering
\caption{The cost model: measured rows, first-order estimates, and open
entries. \textsc{est} marks first-order values from literature figures
(HBM at 3.97~pJ/bit~\cite{oconnor2017fgdram}; \cam cells at
0.1--1~fJ/search~\cite{pagiamtzis2006cam,ni2019fetcam,liu2022evacam}),
at $K=1.3\!\times\!10^{5}$, 2{,}048-bit rows, per token per layer.
\textsc{tbd} entries await array-level (Eva-CAM) and GPU-baseline
measurements.}
\label{tab:costmodel}
{\scriptsize\setlength{\tabcolsep}{3pt}
\begin{tabular}{@{}lccp{0.23\columnwidth}@{}}
\toprule
\textbf{Metric} & \textbf{CAM} & \textbf{GPU} & \textbf{Assumptions} \\
\midrule
Energy / search & \makecell[l]{0.03--0.3\,$\mu$J\\\textsc{est}} &
  \makecell[l]{$\sim$67\,mJ\\\textsc{est}} &
  cell figure spans published range; GPU is DRAM access only \\
Latency / search & \makecell[l]{0.3--4.4\,ns/arr.\\\textsc{est}} &
  \makecell[l]{0.63\,ms\\\textsc{est}} &
  reduction net.\ open; GPU is the batch-1 streaming floor \\
Area / $10^{3}$ rows & \textsc{tbd} & n/a &
  cell type (Table~\ref{tab:camspace}), node \\
Coverage / mm$^{2}$, / W & \textsc{tbd} & \textsc{tbd} &
  behaviors $\times$ layers, at fixed detection rate \\
Traffic / token & $d b L_{\mathrm{mon}}$ & $K d b L_{\mathrm{mon}}$ &
  measured here; Eqs.~(\ref{eq:gputraffic}), (\ref{eq:camtraffic}) \\
Bank capacity & \multicolumn{2}{c}{Fig.~\ref{fig:capacity}} &
  measured here; code widths from Sec.~\ref{sec:Results} \\
\bottomrule
\end{tabular}}
\end{table}

% ===== end 07 - Hardware Analysis.tex =====

%%%%%%%%%%%%%%%%%%%%%%%%%%%%%%%%%%%%%%%%%%%%%%%%%%%%%%%%%%%%%%%%%%%%%%
% ===== 08 - Conclusions.tex =====
%%%%%%%%%%%%%%%%%%%%%%%%%%%%%%%%%%%%%%%%%%%%%%%%%%%%%%%%%%%%%%%%%%%%%%
\section{Conclusions}
\label{sec:Conclusions}
\label{subsec:Limitations}
\label{subsec:Conclusions}
\textbf{Limitations:} Nothing here competes with a deployed
classifier, and no end-to-end jailbreak-blocking experiment was run;
the measured claims concern the fidelity and cost of stored-signature
matching, one component of the always-on stage of
Sec.~\ref{sec:Introduction}. The experiments define deception as an
output that contradicts the known correct output, which the available
datasets control for imperfectly. The grid results classify one
activation per response rather than per token; the readout
reproduction of Sec.~\ref{subsec:Watchlist} is per-token. Layer 33 is
a tuned hyperparameter of one model family at one scale, and all grid
results are from one model. The hardware comparison of
Sec.~\ref{sec:Hardware} is capacity and bandwidth arithmetic on vendor
figures, with no contention term for a monitor co-resident with live
inference, and the energy and latency entries of
Table~\ref{tab:costmodel} are first-order estimates from literature
values rather than our measurements.
\dt[regenerate Fig.~5(a)(b) from result files before submission -- the
plotted values were digitized from the 07-30 deck]{}

Always-on behavioral monitoring of a language model is a library-search
problem, and the variable that decides both its coverage and its cost is
$K$, the number of stored signatures compared against each activation. At
$K=1$, which is what is deployed today, the operation is nearly free and
no special hardware is warranted. Published work gives reasons to expect
$K$ to grow -- non-transfer across behaviors, steering wanting a
direction per behavior per domain per layer, feature dictionaries, and
whole-vocabulary readout, which fixes $K\!\approx\!10^{5}$ per layer by
construction -- and also gives reasons to doubt it, which is why the
question is worth measuring rather than asserting.

One methodological result comes before the fidelity numbers. Accuracy
alone does not establish that a memory-based detector is using its
memory: across five variants of an earlier pipeline, 66--100\% of the
accuracy of the two best-scoring configurations survived randomizing
every stored row, so the trained front end was doing most of the
classifying and the array was close to decoration. Reporting front-end
share alongside \auroc separates the two, and it is what moved this work
onto frozen projections.

What a \cam-resident monitor gives up can then be measured, and we
measure it on the two workloads at opposite ends of the range. For
deception detection at $10^{2}$--$10^{3}$ labeled rows, a 384-bit code
inside a plain array holds in-domain \auroc within 0.009 of
full-precision cosine similarity on the raw 8{,}192 numbers, at 0.920
(\signrp ternary) and 0.919 (thermometer) against 0.928; cross-domain,
the same codes give up 0.065 and 0.127. At a matched 384 bits the
ternary code is ahead of the thermometer code by 0.061 on average, which
is the opposite of what a preference for the cheaper cell would predict
and is the experiment we would run next at matched cell count rather
than matched code length. A two-stage shortlist recovers most of the
cross-domain gap in aggregate, but its fixed $k$ does not travel across
bank sizes, which argues for a threshold criterion over a top-$k$ one in
any array that has to serve banks of different sizes.

For whole-vocabulary readout at $10^{5}$ rows per layer by construction,
where the exact readout is its own ground truth and no labels are needed,
a 3-bit Euclidean code at $\sim$$10^{3}$ storage bits per row -- a
$128\times$ compression of an fp16 row -- reproduces 0.96 of the exact
top-1 selections and 0.81 of the per-row firing decisions. Reporting the
closest row survives compression better than reporting which rows cross a
threshold, by 0.15--0.2 at every code width, which is a matchline-mode
result rather than a cell-type one.

On the cost side, monitoring traffic on a GPU crosses an H100's entire
HBM bandwidth at $K\!\approx\!204$K and a Jetson-class part's at
$K\!\approx\!17$K, while the same library at the measured code widths is
6--\qty{33}{\mega\byte} per layer and fits on a die. A first-order
estimate from literature figures puts the energy gap of the search
itself at five to six orders of magnitude -- roughly 67~mJ per token
per layer to stream the whole-vocabulary bank from HBM against
0.03--0.3~$\mu$J to compare it in place -- of which a factor of 64 is
the measured compression and the rest is the substrate
(Sec.~\ref{subsec:CostModelGap}). A library that is inexpensive to
hold and expensive to move is the condition under which searching it
in place is worth the silicon. The remaining work is short and
specific: Eva-CAM characterization at a stated node and a tuned GPU
baseline to turn the estimates into measurements, and -- upstream of
any hardware
-- establishing the minimum $K$ that adequate behavioral coverage
actually requires.

% ===== end 08 - Conclusions.tex =====

\bibliographystyle{IEEEtran}
\bibliography{refs}

% ---------- review-only appendix (With Comments build only) ----------
\reviewonly{%
\clearpage
\section*{\textcolor{aiviolet}{Length-cut log (not in the submission)}}
\begin{pnote}[aiviolet]
\textbf{Baseline (Sep 16):} v7 tree consolidated into this single
file; 16 pages clean (${\approx}$14.7 body $+$ 1.3 references) against
the DATE limit of 6 pages $+$ 1 for references. The full change history
through v7 lives in the archived driver one level up.

\textbf{Round 1 (Sep 16), structural, per IB's calls:} {\bf (1)}
Sec.~III folded into Sec.~II as the ``Cost model'' subsection (labels
kept, so every Sec.~III cross-reference still resolves). {\bf (2)}
Limitations folded into one paragraph at the head of the Conclusions.
{\bf (3)} Old VI-D (library-size sweep -- digitized data, three known
bugs) cut; Fig.~5(c) is now orphaned and flagged. {\bf (4)} Old VII-F
(the measurement list) cut; its cross-references rewritten, one
sentence of remaining work stays in the Conclusions. {\bf (5)} The
four-domain first-pipeline results paragraph in V-B cut (those numbers
were the undergraduates' pipeline, not the ten-domain runs). {\bf (6)}
Decision rule reworded per IB: the score is the fraction of deceptive
rows among the 31 nearest, and \auroc is computed on it directly.
{\bf (7)} All resolved review notes, the old commented-out IV-C, and
the per-section history comments removed from the source; the open
orange flags remain. No figure or table was changed in this round --
those are being decided one at a time.
\end{pnote}
}

\end{document}
